\section{A fixed-level nonordinary Artin coefficient in the actual inverse tower}

The coefficient field in this section is a sufficiently large finite extension
\(E/\mathbf Q_{269}\), with ring of integers \(\mathcal O\).  Write \(p=269\).
Let \(\rho=\rho_{\theta_\chi}\) be the two-dimensional Artin representation
of the cubic theta series used above.  The same argument allows the
quadratic cyclotomic compensation occurring in the ramified local
normalization.  The assertion concerns a fixed tame level and a fixed
integer \(j\); it does not assert compatible choices as an auxiliary ray
modulus grows.

We use the following primary inputs, with their original conventions.
Pan, \emph{On locally analytic vectors of the completed cohomology of
modular curves}, published DOI \url{https://doi.org/10.1017/fmp.2022.1},
Corollary~6.3.6, \S6.4.2, Theorem~6.4.3, Corollary~6.4.4,
Remark~6.4.5 and Theorem~6.4.8, supplies the completed-cohomology
occurrence and the Kirillov-model action.  The author's arXiv version is
\url{https://arxiv.org/pdf/2008.07099}; the applied passages are on
printed pp.~78--84.  The global twist convention is in \S6.1.1,
printed pp.~72--73, and the standard local-system convention is in
\S4.1.3, printed p.~34.
For the actual coefficient tower we use Kings--Loeffler--Zerbes,
\emph{Rankin--Eisenstein classes and explicit reciprocity laws},
arXiv:1503.02888v5, Definition~2.4.3, Theorem~4.5.1 and
Remark~4.5.3, Definition~5.1.1, and Proposition~7.2.1:
\url{https://discovery.ucl.ac.uk/id/eprint/1497960/1/Zerbes-S_1503.02888v5.pdf}.
The local-system and tower actions are checked against Emerton's
author version, \S2.4, printed pp.~12--13:
\url{https://www.math.uchicago.edu/~emerton/pdffiles/coates.pdf}.
These are published inputs; no new proof of the upstream
\(p\)-adic local Langlands correspondence is asserted.

\begin{proposition}[AC-K33: an actual negative-moment Artin quotient]
\label{AC-K33}
Fix \(j\geq0\).  There is a sufficiently small neat principal tame
level \(M\), prime to \(p\), such that, for \(N=Mp^a\)
with \(a\geq1\), the actual two-level KLZ corestriction tower
\[
 \mathcal I_j(N)=
 \varprojlim_r H^1_{\mathrm{et}}\!
 \left(Y(M,Np^r)_{\overline{\mathbf Q}},
       \operatorname{TSym}^j\mathcal H_{\mathcal O}(1)\right),
 \qquad\text{transition }(\operatorname{pr}_1)_*,
 \tag{K33.1}
\]
has a bounded continuous \(G_{\mathbf Q}\)-equivariant quotient
\[
 \mathcal I_j(N)\otimes_{\mathcal O}E
       \longtwoheadrightarrow \rho^\vee .
 \tag{K33.2}
\]
Here \(\mathcal H_{\mathcal O}=(R^1\pi_*\mathcal O)^\vee\)
is the Tate-module local system, with exactly the convention of
Kings--Loeffler--Zerbes.  On this quotient \(U_p'\) is zero and the
source inverse-diamond operator
\(\Lambda_D(u)=\langle u^{-1}\rangle\), \(u\in\mathbf Z_p^\times\),
acts by \(u^{-j-1}\).

Moreover, (K33.1) is the cohomology of the division-torsor coefficient
\(\Lambda(\mathcal H\langle t_N\rangle)\otimes
\operatorname{TSym}^j\mathcal H(1)\).  It is an actual direct coefficient
summand of the full \(\Lambda^{[j]}(\mathcal H)(1)\) construction.
Thus the quotient is available before ordinary projection, at the
negative source moment \(k'_g-j=-1-j\).  Its fixed-level image admits
an \(\mathcal O\)-lattice.  No nonzero image of a particular
Rankin--Iwasawa class, growing-ray bound on denominators, or
two-dimensional singular-regulator assertion is included.
\end{proposition}

\begin{proof}
\emph{1. The global Artin constituent and its local model.}
Apply Pan's convention to \(V=\rho(1)\).  The completed-cohomology
eigenspace attached to \(V\) is globally \(V(-1)=\rho\)-isotypic,
and its \(\mathrm{GL}_2(\mathbf Q_p)\)-centre acts by
\(\det(V)\varepsilon^{-1}\).  The local representation \(V_p\)
is absolutely irreducible: cubic tame inertia has two distinct
characters interchanged by the unramified quadratic Frobenius.
The residual image of the finite Artin factor has order \(6\),
prime to \(269\).  Its splitting field has Galois group \(S_3\)
and unique nontrivial abelian subfield
\(K=\mathbf Q(\sqrt{-3})\).  This is not the quadratic subfield
\(\mathbf Q(\sqrt{269})\) of \(\mathbf Q(\mu_{269})\).
Consequently restriction to \(G_{\mathbf Q(\mu_p)}\) preserves
residual absolute irreducibility.  The cyclotomic twist and the
possible quadratic compensation are trivial on that subgroup.
Oddness and finite ramification are also preserved.  These are
the actual hypotheses needed for Pan's Theorem~6.4.8.

Pan initially supplies an arbitrary open tame level \(K^p\).
Choose a sufficiently small principal full tame level \(M\),
prime to \(p\), inside \(K^p\).  The finite pullback to this
smaller tame level is injective after tensoring with \(E\):
pushforward composed with pullback is multiplication
by its finite degree.  It commutes with \(G_{\mathbf Q}\) and
the \(p\)-local action.  The resulting curves are the actual
two-level source curves \(Y(M,Np^r)\), \(N=Mp^a\); their
tame part is full principal level \(M\), and their \(p\)-part
is \(\Gamma_1(p^{a+r})\).  We retain the full
\(\mathbf Q\)-scheme, including all Weil-pairing components,
so every pullback and quotient remains \(G_{\mathbf Q}\)-equivariant.
No assertion that an arbitrary
\(K^p\) contains a pure \(\Gamma_1(M)\) level is used.

After this change of tame level and enlargement of \(E\),
Pan's \S6.1.1 gives the global \(\rho\)-isotypic
decomposition.  Corollary~6.3.6 and Schur's lemma let us
choose a nonzero \(\Pi(V_p)\) embedding in
\(\operatorname{Hom}_{G_{\mathbf Q}}(\rho,\widetilde H^1)\).
Choose one such embedding.  All subsequent constructions commute with
the global \(G_{\mathbf Q}\)-action; the chosen multiplicity vector
does not identify that action with a local action.

Put \(W=\mathcal H^\vee=R^1\pi_*\mathcal O\).
Pan's \S4.1.3 identifies \(W\), rather than \(\mathcal H\), with
the standard algebraic two-dimensional representation.  In
particular the dual of \(\operatorname{TSym}^j\mathcal H\) is
\(\operatorname{Sym}^j W\), with its actual standard action.

\emph{2. A vector killed by the source \(U_p\).}
In Pan's Hodge--Tate convention, \(V_p=\rho_p(1)\) has both
weights \(-1\).  At a sufficiently large cyclotomic stage,
\[
 D_{\mathrm{dif}}^+(V_p)
   =tD_{\mathrm{dR}}(V_p)\otimes_E E_\infty[[t]].
 \tag{K33.3}
\]
Choose \(0\ne d\in D_{\mathrm{dR}}(V_p)\).  In the compact
Kirillov model, define the following vector for
\(x\in\mathbf Z_p^\times\), and make it zero otherwise:
\[
 \Phi_{j,d}(x)=
 \sum_{k=0}^j(-1)^k\frac{j!}{(j-k)!}\,
 t^{-k}x^{-k}d\otimes e_1^k e_2^{j-k}.
 \tag{K33.4}
\]
The relation \(\sigma_a(t)=at\) gives the required cyclotomic
equivariance.  Since
\[
 E_{12}(e_1^k e_2^{j-k})
    =(j-k)e_1^{k+1}e_2^{j-k-1},
\]
consecutive terms cancel in
\((tx+E_{12})\Phi_{j,d}\); the sole remaining term is a
multiple of \(td\), which is zero modulo (K33.3).
Hence upper unipotent matrices fix (K33.4).  For
\(\operatorname{diag}(a,1)\), \(a\in\mathbf Z_p^\times\),
the factor \(a^{-k}\) from the Kirillov variable cancels
the factor \(a^k\) from the standard algebraic representation.
Thus (K33.4) is fixed by
\[
 P_0=\left\{
 \begin{pmatrix}a&b\\0&1\end{pmatrix}:
 a\in\mathbf Z_p^\times,\ b\in\mathbf Z_p\right\}.
\]

The source operator is
\(\sum_{i=0}^{p-1}\left(\begin{smallmatrix}p&i\\0&1\end{smallmatrix}\right)\).
For this operator the algebraic \(p^k\) cancels the
\((px)^{-k}\) in (K33.4), while the unipotent polynomial
cancels the corresponding exponential.  If \(px\) is a unit,
the remaining coefficient is the sum of all nontrivial
\(p\)-th roots arising from the additive character, hence zero;
otherwise the support is zero.  This proves \(U_p\Phi_{j,d}=0\).
The \(k=0\) term proves \(\Phi_{j,d}\ne0\).
Pan's Theorem~6.4.3 embeds this compact Kirillov vector into
\(\Pi(V_p)^{\mathrm{la}}\).  Using the multiplicity vector
chosen in step~1 gives a finite-dimensional global subrepresentation
\[
 Z\simeq\rho\ \subset\
 \left(\widetilde H^1\otimes_E
       \operatorname{Sym}^j W_E\right)^{P_0,U_p=0}.
 \tag{K33.5}
\]

\emph{3. Descent through the actual compact-support tower.}
Let \(K_r=1+p^rM_2(\mathbf Z_p)\), and let \(K_{1,r}\) be
the subgroup with lower row congruent to \((0,1)\bmod p^r\).
Then \(K_{1,r}=P_0K_r\).  At the chosen neat tame level,
the map from full \(K_r\)-level to \(K_{1,r}\)-level,
with the fixed principal tame level \(M\), is
a finite etale torsor with group
\(P_r=P_0/(P_0\cap K_r)\).
For finite locally constant coefficients on these affine
curves, \(H_c^0=0\).  Compact-support Cartan--Leray in degree
one therefore identifies quotient cohomology with invariant
cover cohomology.  This uses no division by \(\#P_r\).

For \(r\geq n\), full \(p^r\)-level trivializes
\(\operatorname{Sym}^j W/\!p^n\).  Consequently
\[
 A_n:=\varinjlim_r
 H_c^1\!\left(Y(M,Mp^r)_{\overline{\mathbf Q}},
       \operatorname{Sym}^j W/\!p^n\right)
 =
 \left(\varinjlim_r H_c^1(Y(K_r),\mathcal O/\!p^n)
       \otimes\operatorname{Sym}^j W/\!p^n\right)^{P_0}.
 \tag{K33.6}
\]
Invariants commute with this filtered limit: a finite-level
representative has an open stabilizer; the finitely many
coset differences witnessing invariance vanish at one common
later level.  Thus
\[
 B_0:=\varprojlim_n A_n
   =\left(\widetilde H_{c,\mathcal O}^1
        \otimes\operatorname{Sym}^j W\right)^{P_0}.
 \tag{K33.7}
\]

We verify the completion step instead of assuming that finite-level
\(H_c^2\) is torsion-free.  The direct limit of
\(H_c^2(Y(M,Mp^r),\operatorname{Sym}^j W/\!p^n)\) is zero.
Indeed its pullback transition is dual to trace on \(H^0\) of
the finite dual local system.  For \(r\geq n\), the coset
representatives for the next geometric \(\Gamma_1\)-level
are congruent to \(1\) modulo \(p^r\); their number is \(p^2\).
They act trivially on the coefficient fibre modulo \(p^n\),
so trace is multiplication by \(p^2\).  Repeated transitions
kill the finite coefficient group.

Taking direct limits in the coefficient exact sequences now gives
\[
 0\longrightarrow A_n\longrightarrow A_{n+m}
       \longrightarrow A_m\longrightarrow0
 \tag{K33.8}
\]
with the usual multiplication and reduction maps.
Compatible lifting shows that \(B_0\) is \(p\)-torsion-free
and \(B_0/p^nB_0=A_n\).  This establishes the precise bounded
completion needed below, including its integral lattice.

\emph{4. Duality gives an actual continuous quotient.}
Finite-coefficient Poincare duality gives perfect pairings
\[
 H^1\!\left(Y(M,Mp^r),
       \operatorname{TSym}^j\mathcal H(1)/p^n\right)
 \times
 H_c^1\!\left(Y(M,Mp^r),
       \operatorname{Sym}^jW/p^n\right)
 \longrightarrow\mathcal O/p^n.
 \tag{K33.9}
\]
They are \(G_{\mathbf Q}\)-equivariant and corestriction is
transpose to pullback.  Inverse limits in \(r\) and \(n\)
commute.  The comparison with integral finite-level cohomology
uses the Mittag--Leffler property of the finite \(H^0\)-systems.
Using (K33.8), one obtains the actual bounded dual
\[
 \mathcal I_j(N)
   =\operatorname{Hom}_{\mathcal O,\mathrm{cont}}(B_0,\mathcal O).
 \tag{K33.10}
\]
Changing the initial positive \(p\)-power level is a cofinal
change of tower.

The localized boundary cohomology vanishes at the present
non-Eisenstein ideal.  For completeness, cusps have abelian
Galois action; if the residual Artin ideal occurred there,
Eichler--Shimura at the unramified primes and Chebotarev
would give a reducible two-character residual trace.
Brauer--Nesbitt would contradict the absolute irreducibility
of \(\overline\rho\).  Thus (K33.5) lies in \(B=B_0\otimes E\).

The finite-dimensional subspace \(Z\subset B\) is closed.
Over the finite extension \(E/\mathbf Q_p\), bounded linear
functionals on \(Z\) extend to \(B\) by the nonarchimedean
Hahn--Banach theorem.  Restriction of the bounded dual is
therefore onto \(Z^\vee\).  After scalar multiplication an
extension is integral on \(B_0\), so it comes from (K33.10).
Restriction itself is \(G_{\mathbf Q}\)-equivariant; no
equivariant choice of extensions is required.  It is also
continuous for the compact inverse-tower topology: modulo
\(p^n\), an element of \(B_0\) is represented at a finite
level in (K33.6), so evaluation is a finite-level continuous
operation.  This proves (K33.2) and the fixed-level lattice
assertion.

\emph{5. Division torsors, rather than an unrelated full-level model.}
Here the hypothesis \(p\mid N\) is essential.  The whole
moduli proof of KLZ Theorem~4.5.1, in its stated generality \(M\mid N\)
and \(p\mid N\), identifies
\(E[p^r]\langle t_N\rangle\) over \(Y(M,N)\) with
\(Y(M,Np^r)\), with exactly
the transition \(\operatorname{pr}_1\).
The finite pushforward projection formula, tensoring by
\(\operatorname{TSym}^j\mathcal H\), therefore identifies
(K33.1) with the division-torsor cohomology stated in the
proposition.

Multiplication by \(N=Mp^a\) identifies the inverse division
torsor with the clopen Tate-module coset
\[
 v\bmod p^a=[M]t_N.
 \tag{K33.11}
\]
Its points are primitive modulo \(p\).  Differences form the
Tate module, while multiplication by \(N\) has image
\(N\mathcal H=p^a\mathcal H\); the inverse identification is
the corresponding shift of indices by \(a\).
Restriction of measures to (K33.11) is an integral direct
summand projection from the full measure sheaf.
After tensoring with \(\operatorname{TSym}^j\mathcal H\),
it is precisely a summand of
\(\Lambda^{[j]}(\mathcal H)\).  The resulting cohomology
projection is surjective.  This proves availability of the
coefficient map, not survival of any specified cohomology
class under it.

\emph{6. The inverse-diamond sign and the Hecke transpose.}
Write \(d_a(E,P)=(E,aP)\).
Emerton's actual local-system action is
\[
 (\tau,g,v)\cdot g'=(\tau,gg',g_p'^{-1}v).
 \tag{K33.12}
\]
The Tate frame is \(\alpha_\tau g^{-t}\), so the action
\(D(a)=\operatorname{diag}(1,a)\) on direct compact cohomology
is \(D_B(a)=d_{a^{-1}}^*\).
The unit central character of \(\Pi(V_p)\) is
\(\det(V)\varepsilon^{-1}=a\): the determinant of the
Artin representation, and also that of its quadratic
compensation, is unramified at \(269\).
The algebraic factor adds \(a^j\).  Since \(P_0\) fixes
(K33.4), \(D_B(a)\) acts on (K33.5) by \(a^{j+1}\).

It remains to take the contragredient, not merely the
transpose of \(D_B(a)\).  For the pairing (K33.9),
\[
 \begin{split}
 \langle D_{\mathcal I}(a)m,v\rangle
 &=\langle m,D_B(a^{-1})v\rangle\\
 &=\langle m,d_a^*v\rangle
  =\langle d_{a,*}m,v\rangle.
 \end{split}
 \tag{K33.13}
\]
Thus \(D_{\mathcal I}(a)=d_{a,*}=d_{a^{-1}}^*\), which is
exactly KLZ's operator \(\Lambda_D(a)=\langle a^{-1}\rangle\)
in Proposition~7.2.1.  Its eigencharacter is
\(a^{-j-1}\), not \(a^{j+1}\).
Similarly, the \(p\)-Hecke correspondence on the direct
tower is transpose to \(U_p'\) by KLZ Definition~2.4.3.
The quotient therefore has \(U_p'=0\), as asserted.
\end{proof}


\begin{remark}[The simultaneous two-copy coefficient quotient]
The construction has a useful linear refinement.  Fix the
\(\Pi(V_p)\) embedding in the global multiplicity space
\(\operatorname{Hom}_{G_{\mathbf Q}}(\rho,\widetilde H^1)\)
used in step~1.  Formula (K33.4) is \(E\)-linear in \(d\);
its \(k=0\) coefficient makes this map injective on the
two-dimensional vector space \(D_{\mathrm{dR}}(\rho_p(1))\).
Consequently the same completed coefficient space contains
\[
 \rho\otimes_E D_{\mathrm{dR}}(\rho_p(1))
 \lhook\joinrel\longrightarrow B,
 \tag{K33.15}
\]
where the de Rham space is a multiplicity label with trivial
global Galois action.  Restriction of the bounded dual in
step~4 is simultaneously surjective on this finite-dimensional
subspace.  Therefore the very same actual tower admits
\[
 \mathcal I_j(N)\otimes E
 \longtwoheadrightarrow
 \rho^\vee\otimes_E D_{\mathrm{dR}}(\rho_p(1))^\vee .
 \tag{K33.16}
\]
Both coefficient copies have the same \(U_p'=0\) and
inverse-diamond character.  This supplies two available
coefficient functionals; it does not prove that the images
of the Rankin--Iwasawa class have linearly independent
singular regulators.
\end{remark}
\begin{remark}[Class survival and the two distinct towers]
The ordinary class of KLZ Theorem~6.3.4 is obtained by
ordinary projection followed by powers of
\((U_p',U_p')^{-1}\); its projection to (K33.2) is zero.
This does not settle the preordinary Rankin--Iwasawa image.

More precisely, KLZ Theorem~5.3.1 gives compatibility for
the internal \(\operatorname{pr}_1\)-level tower.
The distinct cyclotomic norm relation of Theorem~5.4.1
and the definition in \S6.1 give, after any compatible
fixed-level coefficient projection with
\(A=U'_{p,f}U'_{p,g}=0\),
\[
 \operatorname{cor}(z_{mp^{r+1}})=0\quad(r\geq1),
 \qquad
 \operatorname{cor}(z_{mp})=-p^j\sigma_p z_m
 \quad(p\nmid m).
 \tag{K33.14}
\]
The first-step term in (K33.14) is nonzero as an operator
on a nonzero characteristic-zero coefficient space.
These identities alone do not prove that \(z_m\) is zero
or nonzero.  The actual image, uniform auxiliary-ray
compatibility and the rank of the resulting local singular
regulator remain separate gates.
\end{remark}

\begin{remark}[Geometric dictionary and precise scope]
The division torsors in step~5 are genuine finite etale
covers and their limit is the actual measure coefficient
used by the Rankin--Iwasawa construction.  The quotient
(K33.2) is a continuous cohomological coefficient map.
It is not an assertion that a negative algebraic symmetric
power is a pure motive, and it does not contradict the
geometric de Rham obstruction proved in AC-G69.
The vector in step~2 is locally analytic and is used
before ordinary projection.  No Grothendieck archive
locator or additional motivic realization is asserted.
\end{remark}
